\documentclass[utf8]{frontiers_suppmat}
\usepackage{url,hyperref,lineno,microtype,amsmath,booktabs}
\usepackage[onehalfspacing]{setspace}

\begin{document}
\onecolumn
\firstpage{1}

\title[Supplementary Material]{{\helveticaitalic{Supplementary Material: SIMPLEX Composition-Space Deep Learning for Hydrogel Adhesion}}}

\maketitle

\section{Supplementary Data}

This supplement provides the complete experimental and analytical details
behind the SIMPLEX framework summarised in the main article: dataset
description and the training/prospective split, architecture configuration,
training protocol, the full ablation table, prospective (external) screening
statistics on the 25-formulation held-out cohort, multi-target extensibility,
interpretation highlights, and a reproduction guide. All code, processed
tensors, per-formulation predictions and regeneration scripts are available
at \url{https://github.com/SHENTongfei/HydroGelNet}.

\section{Data and Preprocessing}

\subsection{Raw dataset}
The raw experimental data were obtained from the public repository
\emph{sheng-hu/hydrogels} (MIT licence) accompanying Liao et al.,
\textit{Nature} 644, 89--95 (2025), DOI 10.1038/s41586-025-09269-4. The
dataset records underwater glass adhesion strength (kPa) for 341 bio-inspired
hydrogel formulations, each described by the molar fractions of six
functional monomers:
\begin{itemize}
\item 2-hydroxyethyl acrylate (HEA, nucleophilic);
\item butyl acrylate (BA, hydrophobic);
\item 2-carboxyethyl acrylate (CBEA, acidic);
\item (3-acrylamidopropyl)trimethylammonium chloride (ATAC, cationic);
\item 2-phenoxyethyl acrylate (PEA, aromatic);
\item acrylamide (AAm, amide).
\end{itemize}
The six fractions sum to 1 (composition simplex).

\subsection{Training / prospective split}
\begin{itemize}
\item \textbf{Training (internal cohort):} the 316-formulation set
  ($df\_316$). Adhesion spans the full experimentally explored range,
  1~--~353~kPa (mean 93~kPa), with 19\% of samples above 150~kPa, so the
  model is trained on the complete composition--adhesion surface rather than
  on a low-adhesion subset.
\item \textbf{Prospective (external cohort):} the 25 formulations added in
  the final round of the original study's sequential model-based optimisation
  (SMBO) loop, obtained by matching compositions between $df\_341$ and
  $df\_316$ ($df\_341$ minus $df\_316$; adhesion 62~--~251~kPa, mean
  158~kPa). These are the last model-guided discoveries and were never
  touched during model selection, hyper-parameter tuning, or ablation.
\end{itemize}
The prospective cohort is a \emph{model-guided, genuinely held-out} cohort
rather than an independent-laboratory cohort: it originates from the same
laboratory and instrumentation (the final SMBO round). This is acknowledged
as a limitation; cross-laboratory validation remains an important next step.
The cohort is small ($n=25$), so per-metric estimates carry uncertainty and
are reported with bootstrap confidence intervals (Section~5).

\subsection{Features}
Two modalities, 21 features total: (1) \emph{monomer fractions} (6 raw molar
fractions); (2) \emph{pairwise synergy} (15 products $x_i x_j$, $i<j$),
encoding composition synergy explicitly following compositional-data and
materials-informatics practice. Scaling and imputation were fitted inside
each cross-validation training fold only (no leakage); the prospective cohort
was standardised with statistics from the training cohort only.

\section{Architecture and Training Protocol}

SIMPLEX embeds the two modalities to a common dimension ($d=64$), refines the
fused representation with two residual blocks separated by an interaction
self-attention layer (4 heads), pools, and maps through a linear head to a
non-negative adhesion prediction under a range-domain constraint. The final
configuration (Table~\ref{tab:hyper}) was selected using internal CV only;
every switch was ablation-gated, and any switch that did not pay for itself
was removed from the final model.

\begin{table}[h]
\centering
\caption{Final SIMPLEX hyper-parameters (best\_config.json; searched and frozen on the internal cohort only).}
\label{tab:hyper}
\begin{tabular}{ll}
\toprule
Hyper-parameter & Value \\
\midrule
Embedding dimension $d$ & 64 \\
Residual blocks & 2 \\
Interaction self-attention heads & 4 \\
Dropout & 0.15 \\
Fusion & concat (dual-modality, interaction-aware) \\
Backbone & residual blocks (Transformer: off, ablation-gated) \\
FiLM / modality gate & off \\
Mixup & on, $\alpha = 0.4$ \\
Stochastic weight averaging (SWA) & on (start 0.6, lr fraction 0.25) \\
Range-domain constraint & on ($w = 0.1$) \\
SAM / EMA / R-Drop / feature noise & off \\
Contrastive / MFM / reconstruction pre-training & off \\
Uncertainty weighting & off (single target) \\
Target transform / scaler & standard / standard \\
Max epochs / patience & 200 / 30 \\
Batch size & 32 \\
Learning rate & $3\times10^{-3}$ \\
Weight decay & $10^{-3}$ \\
Gradient clipping & 1.0 \\
\bottomrule
\end{tabular}
\end{table}

Design notes: (i) \textbf{Residual blocks, not a Transformer backbone} ---
on this small-sample regime the Transformer backbone did not pay for itself
and was removed by the ablation gate. (ii) \textbf{Mixup is the key
regulariser} --- removing it degrades internal $R^2$ by $+0.067$ (Section~4).
(iii) \textbf{Self-supervised pre-training (contrastive/MFM/reconstruction)
was neutral or detrimental} and was removed (ablation-gated). (iv) The
range-domain constraint and Mixup keep predictions physically plausible
(non-negative, in-range) on the prospective cohort.

\section{Full Ablation Table}

Leave-one-out ablation was run under the search protocol (2 seeds $\times$ 3
folds); the full-model mean internal $R^2$ in this protocol is 0.713, which
differs from the 5-seed main-CV mean (0.793, Section~2/main article) because
the cheaper search protocol uses fewer seeds. Table~\ref{tab:abl} reports the
per-variant deltas; a positive delta means removing the component degrades
performance (component helps).

\begin{table}[h]
\centering
\caption{Full ablation (internal $R^2$; $\Delta$ = full model $-$ variant; full mean 0.713, 2 seeds $\times$ 3 folds).}
\label{tab:abl}
\begin{tabular}{lc}
\toprule
Variant & $\Delta R^2$ (full $-$ variant) \\
\midrule
w/o multimodal fusion & $+0.093$ (largest) \\
w/o Mixup & $+0.067$ \\
fusion = gated & $+0.062$ \\
w/o residual blocks & $+0.062$ \\
fusion = film & $+0.030$ \\
w/o task-specific gating & $+0.022$ \\
fusion = cross & $+0.004$ \\
w/o attention sparsity reg. & $-0.000$ \\
w/o sparse attention & $-0.008$ \\
\bottomrule
\end{tabular}
\end{table}

Fourteen further candidate mechanisms had exactly zero delta in this
protocol --- EMA, FiLM conditioning, MC-Dropout, MFM pre-training, R-Drop,
SAM, SWA, contrastive pre-training, feature noise, modality gate, pretrained
transfer, transformer block, uncertainty weighting, and the range-domain
constraint --- and the two attention-related switches above (attention
sparsity regularisation, sparse attention) had negative delta. In total,
sixteen candidate switches showed zero or negative contribution. Consistent
with the ablation-gating principle, switches with non-positive contribution
were not adopted as primary components of the final model
(Table~\ref{tab:hyper}); the range-domain constraint is nevertheless retained
at zero measurable cost because it enforces physically plausible, in-range
predictions on the prospective cohort. The ablation spread across variants is
informative and the retained switches are reported transparently.

\section{External Prospective Statistics (25 held-out formulations)}

The 25-formulation prospective cohort was evaluated \emph{once}, with the
frozen ensemble, after all hyper-parameters had been selected. Metrics follow
the screening application: absolute $R^2$, Spearman rank correlation,
tier-level ROC-AUC (top half vs bottom half of the cohort), and Top-$k$
screening precision (fraction of the true top-$k$ formulations recovered
among the model's top-$k$ predictions).

\subsection{Point estimates and bootstrap confidence intervals}

Sample-level cluster bootstrap (25 formulations, 2000 resamples;
bootstrap\_ci.csv):

\begin{table}[h]
\centering
\caption{SIMPLEX prospective performance with 95\% bootstrap CIs.}
\label{tab:ext_ci}
\begin{tabular}{lccc}
\toprule
Metric & Point & 95\% CI & SE \\
\midrule
$R^2$ & 0.710 & [0.458, 0.857] & 0.104 \\
RMSE (kPa) & 33.66 & [24.75, 41.40] & 4.33 \\
MAE (kPa) & 26.88 & [19.52, 34.79] & 4.03 \\
\bottomrule
\end{tabular}
\end{table}

On the same cohort, SIMPLEX attains Pearson $r = 0.877$, Spearman
$\rho = 0.87$, tier-level ROC-AUC $= 0.94$, Top-10 precision $= 1.00$ and
Top-20 precision $= 0.95$: all ten highest-predicted formulations are among
the ten strongest adhesives in the held-out cohort.

\subsection{Model comparison}

\begin{table}[h]
\centering
\caption{Prospective screening metrics on the 25 final model-discovered formulations.}
\label{tab:ext}
\begin{tabular}{lccccc}
\toprule
Model & $R^2$ & Spearman $\rho$ & ROC-AUC & Top-10 & Top-20 \\
\midrule
\textbf{SIMPLEX} & \textbf{0.71} & \textbf{0.87} & \textbf{0.94} & \textbf{1.00} & \textbf{0.95} \\
SVR-RBF & 0.71 & 0.85 & 0.94 & 1.00 & 0.90 \\
Ridge & 0.64 & 0.81 & 0.93 & 1.00 & 0.90 \\
ElasticNet & 0.63 & 0.85 & 0.93 & 1.00 & 0.90 \\
GBR & 0.52 & 0.80 & 0.94 & 0.90 & 0.90 \\
RandomForest & 0.44 & 0.80 & 0.92 & 1.00 & 0.90 \\
\bottomrule
\end{tabular}
\end{table}

SIMPLEX attains the highest prospective $R^2$ (tied with SVR, above Ridge and
random forest), the highest Spearman correlation, and the highest Top-20
precision of all models. The deep model's advantage over tree ensembles is
structural: random forest, despite matching SIMPLEX internally ($0.81$ vs
$0.79$), drops to $R^2 = 0.44$ on the prospective cohort because its
nearest-neighbour-style extrapolation fails outside the density of the
training composition space, whereas SIMPLEX's continuous, interaction-aware
composition encoding transfers with a small internal-to-prospective
generalisation gap ($0.79 \rightarrow 0.71$).

\section{Multi-target Extensibility}

The original dataset also records rheological targets for the training
formulations: storage-related modulus (kPa) and loss tangent tan-delta.
$\mathrm{corr}(\mathrm{Glass}, \mathrm{Modulus}) = -0.036$ --- independent
information not used by the Nature study's single-target models. A
single-task SIMPLEX trained on Modulus achieves internal CV $R^2 = 0.396$ vs
RandomForest 0.392 (3 seeds), i.e.\ parity with the strongest baseline on a
second, independent target. Joint multi-target training is left for future
work (see Limitations).

\section{Interpretation Highlights}

Cross-validated permutation importance (mean over folds/seeds; 5 folds
$\times$ 2 seeds) ranks the monomer and interaction features
(Table~\ref{tab:interp}); the direction column reports the sign of the
feature--target association from the candidate-marker screen.

\begin{table}[h]
\centering
\caption{Permutation importance (mean $\pm$ SE) and association direction.}
\label{tab:interp}
\begin{tabular}{clcc}
\toprule
Rank & Feature & Importance $\pm$ SE & Direction \\
\midrule
1 & pair (BA$\times$PEA), hydrophobic--aromatic & 0.143 $\pm$ 0.036 & $+$ \\
2 & Cationic-ATAC & 0.061 $\pm$ 0.016 & $+$ \\
3 & Hydrophobic-BA & 0.050 $\pm$ 0.013 & $+$ \\
4 & pair (HEA$\times$ATAC) & 0.050 $\pm$ 0.012 & $-$ \\
5 & pair (HEA$\times$BA) & 0.039 $\pm$ 0.015 & $-$ \\
6 & pair (BA$\times$ATAC) & 0.035 $\pm$ 0.012 & $+$ \\
7 & pair (BA$\times$CBEA) & 0.028 $\pm$ 0.008 & $-$ \\
\bottomrule
\end{tabular}
\end{table}

Three reproducible, literature-consistent findings emerge. First, the
BA$\times$PEA (hydrophobic--aromatic) interaction term is the single most
important feature (importance 0.143, FDR-corrected $p < 10^{-53}$), followed
by the cationic monomer ATAC (0.061) and the hydrophobic monomer BA (0.050)
--- consistent with the electrostatic/hydrophobic synergy known to drive
underwater adhesion, and with the model-discovered high-performance region
enriched in the hydrophobic--aromatic (BA--PEA) combination. Second, the
interaction terms involving the hydrophilic monomer HEA (HEA$\times$ATAC,
HEA$\times$BA) and the acidic monomer CBEA (BA$\times$CBEA) are significantly
\emph{negative} (inhibitory), indicating that hydrogen-bond-donating
hydrophilic groups and acidic residues weaken wet adhesion to the glass
substrate --- a hydration-layer screening effect consistent with reported
underwater-adhesion chemistry. Third, the partial-dependence profile of the
top feature (BA$\times$PEA) shows adhesion rising with the BA$\times$PEA
fraction over the explored range, providing a concrete design rule for the
composition simplex. These are statistical associations, not causal
mechanisms; they are reported as hypothesis-generating markers for subsequent
experimental validation.

\section{Reproduction Guide}

\subsection{Environment}
Python 3.11 (conda env \texttt{HydroGelNet}), PyTorch 2.14.0+cu130, NVIDIA
RTX 5080 16~GB. Fixed seeds: [42, 2024, 7, 1337, 20260731]; all preprocessing
is fitted inside CV training folds only.

\subsection{Pipeline}
Full pipeline:
\texttt{download $\rightarrow$ build $\rightarrow$ qc $\rightarrow$ tune
$\rightarrow$ train $\rightarrow$ baselines $\rightarrow$ stats
$\rightarrow$ gate $\rightarrow$ interpret $\rightarrow$ figures
$\rightarrow$ tables $\rightarrow$ paper} (run\_all.py). The train stage
trains the final configuration with 5 seeds and the baselines receive an
identical hyper-parameter tuning budget (40 iterations); stats computes the
Nadeau--Bengio corrected resampled t-test (Holm-adjusted), cluster-bootstrap
CIs, permutation tests and Top-$k$ screening statistics; interpret runs
permutation importance and candidate-marker screens.

\subsection{PERF-GATE verdict}
PASS: G1 internal tie with the best baseline (random forest, $0.79$ vs
$0.81$, Holm-corrected $p = 1.0$), G2 no significant disadvantage, G3 seed
stability, G4 prospective $R^2 = 0.710$ above the best baseline (Ridge
$0.637$), G5 informative ablation.

\subsection{Data and Code Availability}
Data: \emph{sheng-hu/hydrogels} (MIT), Nature 2025 supplement; the
training/prospective split is obtained by matching compositions ($df\_341$
minus $df\_316$ = 25 prospective formulations). Code:
\url{https://github.com/SHENTongfei/HydroGelNet} (to be released).
Reproducible artefacts: \texttt{data/}, \texttt{results/},
\texttt{figures/}, \texttt{tables/}, \texttt{paper/}.

\end{document}
